\documentclass[10pt,a4paper]{amsart}
\usepackage[margin=19mm]{geometry}
\usepackage{amsmath,amssymb,mathtools}
\usepackage[scr=rsfs,cal=euler]{mathalfa}
\usepackage{xfrac}
\usepackage[unicode,hidelinks,bookmarks=false]{hyperref}
\newcommand{\ie}{i.e.\ }
\newcommand{\cf}{cf.\ }
\newcommand{\resp}{respectively\ }
\newcommand{\Subsection}{\subsection}
\providecommand{\nobreakdash}{\nobreak\hbox{-}}
\setlength{\emergencystretch}{2em}
\multlinegap0pt
\usepackage{amssymb}
\usepackage{xcolor}
\newtheorem{remark}{Remark}
\newtheorem{proposition}{Proposition}
\title{Irreducibility of the Cuboid Polynomial $P_{a,u}(t)$ via a Rank-Zero Elliptic Curve}
\author{Valery Asiryan}
\date{Source: arXiv:2510.11768v3 (CC BY 4.0)}
\begin{document}
\maketitle

\noindent\emph{Source and licence.} Irreducibility of the Cuboid Polynomial $P_{a,u}(t)$ via a Rank-Zero Elliptic Curve. Version arXiv:2510.11768v3 (CC BY 4.0), distributed by arXiv under the Creative Commons Attribution 4.0 International licence, http://creativecommons.org/licenses/by/4.0/, as recorded in the arXiv licence metadata for this exact version and shown on the abstract page https://arxiv.org/abs/2510.11768v3. Full licence notice: \texttt{licenses/arxiv-2510.11768-CC-BY-4.0.txt}.\par\smallskip

\noindent\textit{Editorial scope.} Counted unit: the article's proposition prop:Cpoints (source0.tex lines 359--364). The article's complete original proof of it is delivered with it (source0.tex lines 366--383, one \texttt{proof} environment of 18 lines). No mathematics is written, repaired or re-derived: the statement and the proof are the article's own lines, in the article's own notation, and no retained line is substituted or re-flowed.  Retained uncounted as the source-local context the proof needs: lines 210--212; lines 229--243; lines 246--253; lines 309--325; lines 322--324.  Nothing was omitted from any retained span, so the package carries no omission record.  No retained line contains a citation, so the wrapper carries no bibliography and no bibliography entry.  No retained line contains a figure environment, an image-inclusion command or drawing code, so no image is delivered and the package declares no asset dependency.  Editorial formatting only, in addition: an AMS \texttt{amsart} preamble that restates the article's own declarations and macros used by the retained lines, namely the environments \texttt{remark}, \texttt{proposition} on separate counters, together with the standard packages those lines need.  The retained environments print in this document as Remark 1, Proposition 1 in order of appearance, whereas the article prints the same items with the numbering of its own full document.  The complete native source of the article as distributed in the arXiv e-print for this exact version is bundled unchanged as \texttt{sources/186986/source0.tex}.
\begin{equation}\label{eq:E}
E:\quad Y^2=X(X-8)(X-9)=X^3-17X^2+72X.
\end{equation}

\begin{remark}[Torsion on $E$]\label{rem:tors}
For $E$ in \eqref{eq:E}, the three points $(0,0),(8,0),(9,0)$ are rational $2$-torsion.
Moreover, the points $(6,\pm 6)$ and $(12,\pm 12)$ have order $4$.
Indeed, for $E:\ Y^2=X^3-17X^2+72X$ the doubling slope is
\[
m=\frac{3X^2-34X+72}{2Y},
\]
and the duplication formulas are
\[
X(2P)=m^2+17-2X,\qquad Y(2P)=-Y+m\bigl(X-X(2P)\bigr).
\]
For $P=(6,6)$ one gets $m=-2$ and hence $2P=(9,0)$; for $P=(12,12)$ one gets $m=4$ and hence $2P=(9,0)$.
Since $(9,0)$ is a nontrivial $2$-torsion point, these points have order $4$.
Thus $E(\mathbb{Q})$ contains a subgroup isomorphic to $\mathbb{Z}/2\mathbb{Z}\oplus\mathbb{Z}/4\mathbb{Z}$; proposition~\ref{prop:rank0} shows this is the full torsion subgroup.
\end{remark}

\begin{proposition}[Rank $0$ and torsion]\label{prop:rank0}
For the curve $E/\mathbb{Q}$ given by \eqref{eq:E}, one has
\[
\mathrm{rank}\,E(\mathbb{Q})=0,\qquad
E(\mathbb{Q})_{\mathrm{tors}}\cong \mathbb{Z}/2\mathbb{Z}\oplus\mathbb{Z}/4\mathbb{Z},\qquad
\mathrm{Cond}(E)=48.
\]
\end{proposition}

\begin{proposition}[Explicit isomorphism $\overline{\mathcal{C}}\simeq E$]\label{prop:CEiso}
Define a rational map $\phi:\mathcal{C}\dashrightarrow E$ on the affine chart by
\begin{equation}\label{eq:phi}
\phi(y,v)=(X,Y),\qquad
X=\frac{v+4y^2+17}{2},\qquad
Y=y\,(v+4y^2+17).
\end{equation}
Then $(X,Y)$ satisfies $E:\ Y^2=X(X-8)(X-9)$. Moreover, $\phi$ extends to an isomorphism of smooth projective curves $\overline{\mathcal{C}}\to E$ over $\mathbb{Q}$ sending
\[
P_\infty^+=(1:4:0)\longmapsto O,\qquad
P_\infty^-=(1:-4:0)\longmapsto (0,0).
\]
The inverse map on affine points with $X\neq0$ is
\begin{equation}\label{eq:phi_inv}
y=\frac{Y}{2X},\qquad v=X-\frac{72}{X}.
\end{equation}
\end{proposition}

\begin{equation}
y=\frac{Y}{2X},\qquad v=X-\frac{72}{X}.
\end{equation}

\begin{proposition}\label{prop:Cpoints}
The set $\overline{\mathcal{C}}(\mathbb{Q})$ consists of exactly $8$ points: the two points at infinity $P_\infty^\pm$ and the six affine points
\[
(y,v)=(0,\pm1),\qquad \Bigl(\pm\tfrac12,\ \pm6\Bigr).
\]
\end{proposition}

\begin{proof}
By proposition~\ref{prop:rank0} and remark~\ref{rem:tors}, the Mordell--Weil group is
$E(\mathbb{Q})=E(\mathbb{Q})_{\mathrm{tors}}$ of order $8$, and its rational points are
\[
E(\mathbb{Q})=\{O,\ (0,0),\ (8,0),\ (9,0),\ (6,\pm6),\ (12,\pm12)\}.
\]
Via the explicit isomorphism $\overline{\mathcal{C}}\simeq E$ from proposition~\ref{prop:CEiso} we obtain:
\[
O\longleftrightarrow P_\infty^+,\qquad (0,0)\longleftrightarrow P_\infty^-,
\]
and for $X\neq0$ the inverse formulas \eqref{eq:phi_inv} give
\[
(8,0)\mapsto (0,-1),\quad (9,0)\mapsto(0,1),\quad
(6,\pm6)\mapsto\bigl(\pm\tfrac12,-6\bigr),\quad
(12,\pm12)\mapsto\bigl(\pm\tfrac12,6\bigr).
\]
Thus $\overline{\mathcal{C}}(\mathbb{Q})$ contains precisely the $8$ points listed.
\end{proof}

\end{document}
